\subsection{The improvement problem and popular tops}\label{forest:section}

We now turn to the improvement estimate \eqref{red:model-improved}
in the local Lipschitz setting: \(|u|\le1\),
\(\Lip(u)\le C_{\rm Lip}\), and the sets \(F,E\) have positive
finite measure and lie in the horizontal testing strip
\(\mathcal D=I\times\R\), where \(|I|\le b_0\).
Write \(\mathcal H\) for the finite-dimensional Hilbert space of
vertical bins; the input and test satisfy
\(\norm{f(z)}_{\mathcal H}\le\one_F(z)\) and
\(\norm{g(z)}_{\mathcal H}\le\one_E(z)\).
We use the parameters in \eqref{red:small-time-parameters}: \(L\)
is sufficiently large, \(m\asymp\exp(L^{0.9})\), \(W=mw\),
and every tile length is at most the dyadic cutoff
\(\tau\le\exp(-L^{0.97})\). All tile collections in this part
are finite. Every estimate is independent of the number of tiles,
tops, scales, and bins.

\subsection{The assignment to popular tops}

Using the fixed exponents \(b_4=1\) and \(c_4=1/4\) from
\eqref{tile:eq-improvement-loss}, choose once an integer \(C_2\)
with \(c_4C_2\ge b_4+42\). Remove tiles
for which \(D_E(s)\le L^{-C_2}\). By
\Cref{tile:prop-model}, in the form
\eqref{tile:eq-improvement-loss}, the absolute
value of their contribution is at most
\[
 C L^{b_4-c_4C_2}\norm{f}_2\sqrt{|E|}
 \le C L^{-40}\sqrt{|F||E|}.
\]
The choice of \(C_2\) depends only on the fixed ordinary loss and
the fixed target exponent.

Apply \Cref{tile:lem-density} to the remaining tiles with threshold
\(\eta=L^{-C_2}\), retaining its unsplit keeper boxes. Denote
the resulting finite collection of tops by \(\mathcal T\).
A top \(T\) has center \(z_T=(x_T,y_T)\), slope \(s_T\),
length \(l_T\le\tau\), and box
\[
 R_T=\left\{(x,y): |x-x_T|\le\frac{l_T}{2},\quad
 |(y-y_T)-s_T(x-x_T)|\le\frac W2\right\}.
\]
We put \(d_T=w/l_T\), and denote the top assigned to a tile \(s\)
by \(T(s)\). Set \(C_{\mathrm{as}}=C_2+10\), increasing the
absolute lower threshold for \(L\) if necessary. This exponent
absorbs the assignment dilation
\(C(1+L^{C_2/100})\), the witness density
\(cL^{-C_2}\), and the area bound \(CL^{C_2}|E|\).
The density lemma gives
\begin{equation}\label{forest:eq-assignment}
 \ell_s\le l_{T(s)},\qquad s_{T(s)}\in\Omega_s,\qquad
 |\theta_s-s_{T(s)}|\le C d_s,\qquad
 z_s\in L^{C_{\mathrm{as}}}R_{T(s)}.
\end{equation}
There are pairwise disjoint measurable sets \(E_T\) such that
\begin{equation}\label{forest:eq-witnesses}
 \begin{split}
 &E_T\subset E\cap L^{C_{\mathrm{as}}}R_T,
 \qquad |E_T|\ge L^{-C_{\mathrm{as}}}|R_T|,\\
 &|u(z)-s_T|\le C d_T\quad(z\in E_T),\qquad
 S:=\sum_{T\in\mathcal T}|R_T|
       \le L^{C_{\mathrm{as}}}|E|.
 \end{split}
\end{equation}
These inequalities hold for every sufficiently large \(L\),
without additional constants in the displayed powers. The angular inclusion in
\eqref{forest:eq-assignment} follows from the nesting of the keeper
interval inside \(\Omega_s\). The direction bound on \(E_T\)
follows from the length of that keeper interval. The disjointness and
lower measure bound also imply the stated estimate for \(S\).

The horizontal centers of all tops lie in a fixed slight enlargement
of \(I\). Indeed, \(E_T\ne\varnothing\), and a point of
\(E_T\) has horizontal coordinate in \(I\) and at distance at
most \(L^{C_{\mathrm{as}}}l_T/2\) from \(x_T\). Since
\(L^{C_{\mathrm{as}}}\tau\to0\), this distance is smaller
than any prescribed positive absolute constant for sufficiently large
\(L\).

\subsection{The adjoint sum and exceptional subsets}

Only the retained tiles, with the assignment \(T(s)\), are used
henceforth. With pairings linear in the first entry, define
\begin{equation}\label{forest:eq-adjoint}
 a_{s,j}=|R_s|^{-1}
       \langle g_j,b_{s,j}(u)\Psi_{s,j}\rangle,
 \qquad \phi_s=(a_{s,j}\Phi_{s,j})_j,
 \qquad P=\sum_s\phi_s.
\end{equation}
These conventions give precisely the form \eqref{eq:src-3}:
\[
 \langle f,P\rangle
 =\sum_s|R_s|^{-1}\sum_j
       \langle f_j,\Phi_{s,j}\rangle
       \langle b_{s,j}(u)\Psi_{s,j},g_j\rangle.
\]
Indeed \(\langle f_j,a_{s,j}\Phi_{s,j}\rangle
=\overline{a_{s,j}}\langle f_j,\Phi_{s,j}\rangle\).
For fixed \(g\) satisfying
\(\norm{g(z)}_{\mathcal H}\le\one_E(z)\), the supremum of
the absolute value of this pairing over
\(\norm{f(z)}_{\mathcal H}\le\one_F(z)\) equals
\(\int_F\norm{P(z)}_{\mathcal H}\dd z\). Equality follows
by taking \(f=P/\norm{P}_{\mathcal H}\) on
\(F\cap\{P\ne0\}\), and zero elsewhere.

Put \(N=\sqrt{|E|/|F|}\). The remaining objective, to be proved
by the subsequent estimates, is
\begin{equation}\label{eq:src-7}
 \int_F\norm{P(z)}_{\mathcal H}\dd z
       \le C L^{-20}N|F|,
 \qquad \frac L2\le\log N\le2L.
\end{equation}
Applying \eqref{tile:eq-improvement-loss} with its absolute density upper bound
and arbitrary \(L^2\) input, and then taking the Hilbert-space
dual norm, gives the fixed exponent \(C_P=b_4=1\) such that
\begin{equation}\label{forest:eq-adjoint-l2}
 \norm{P}_{L^2(\R^2;\mathcal H)}
       \le C L^{C_P}\sqrt{|E|}.
\end{equation}
Consequently, for every measurable \(Z\subset F\),
\begin{equation}\label{forest:eq-discard}
 \int_Z\norm{P(z)}_{\mathcal H}\dd z
 \le C L^{C_P}
       \left(\frac{|Z|}{|F|}\right)^{1/2}N|F|.
\end{equation}
For any prescribed fixed \(Q_0>0\), an exceptional subset with
\(|Z|\le L^{-2(C_P+Q_0+1)}|F|\) therefore costs at most
\(C L^{-Q_0}N|F|\). This estimate is applied to the original
sum \(P\) before any masks or other modifications. It does not
require an \(L^2\) bound for those modified sums.

\subsection{A weak \texorpdfstring{\(L^2\)}{L2} estimate for the top count}

For a top \(T\), use normalized coordinates
\begin{equation}\label{forest:eq-coordinates}
 X_T(z)=\frac{x-x_T}{l_T},\qquad
 Y_T(z)=\frac{(y-y_T)-s_T(x-x_T)}{W},\qquad
 \dist_T(z)=|X_T(z)|+|Y_T(z)|.
\end{equation}
Fix a sufficiently large absolute exponent \(D_2>4\). For now let
\(K_0\ge1\) be any fixed power of \(L\); the choice used with
the masks is specified below. Set \(C_3=2^{D_2}\) and
\begin{equation}\label{eq:src-5}
 h_T(z)=C_3\left(1+\frac{\dist_T(z)}{K_0}\right)^{-D_2},
 \qquad n(z)=\sum_{T\in\mathcal T}h_T(z).
\end{equation}
In particular, \(h_T\ge1\) on \(\{\dist_T\le K_0\}\).
We use the convention
\[
 \norm{a}_{L^{2,\infty}(\mathcal D)}
 =\sup_{\lambda>0}\lambda
   |\{z\in\mathcal D:|a(z)|>\lambda\}|^{1/2}.
\]

\begin{proposition}\label{forest:prop-count}
For the tops in \eqref{forest:eq-witnesses},
\begin{equation}\label{eq:src-6}
 \norm{n}_{L^{2,\infty}(\mathcal D)}
       \le C L^C\sqrt S.
\end{equation}
More precisely, if
\[
 D_*\ge\max(2K_0,\,2L^{C_{\mathrm{as}}},\,2),
\]
then one may take the right-hand side to be
\(C L^{C_{\mathrm{as}}}D_*^3\sqrt S\), provided \(D_*\)
is chosen within an absolute factor of the displayed maximum.
In particular, when \(K_0=L^k\), a permissible exponent in
\eqref{eq:src-6} is
\begin{equation}\label{forest:eq-count-exponent}
 c_{\mathrm{cnt}}(k)
     =C_{\mathrm{as}}+3\max(k,C_{\mathrm{as}}).
\end{equation}
\end{proposition}

Once the fixed power \(K_0\) is chosen, this estimate together with
\(S\le L^{C_{\mathrm{as}}}N^2|F|\) makes the part of \(F\) where
\(n\) exceeds \(N\) times a sufficiently large fixed power of \(L\)
small enough for \eqref{forest:eq-discard}.

\begin{proof}
If \(S=0\), the assertion is immediate, so assume \(S>0\).
We first state the maximal-function input in the form used here.
For a Lipschitz unit vector field \(v_0\), let
\(M_{v_0,\delta}\) be the supremum of averages over Euclidean
rectangles containing the evaluation point, with long side at most
\(c_{\mathrm{LL}}/\Lip(v_0)\), for which the field belongs to
the rectangle's direction uncertainty arc on at least a
\(\delta\) fraction of the rectangle. The arc is centered at the
long-side direction and has total length equal to the short-to-long
side ratio, as in \Cref{red:kakeya}. \Cref{red:kakeya}, namely \cite[Theorem~1.4]{LLKakeya}, gives
\begin{equation}\label{forest:eq-kakeya-input}
 \norm{M_{v_0,\delta}a}_{2,\infty}
       \le C\delta^{-1/2}\norm{a}_2,
       \qquad 0<\delta<1.
\end{equation}
The same estimate holds if the admissible rectangles are required to
satisfy a smaller fixed uncertainty interval, since this restricts the
family of rectangles. When \(\Lip(v_0)=0\), the length restriction
is understood as absent. We apply this result to
\(v_0=(1,u)/\sqrt{1+u^2}\), whose Lipschitz constant is bounded
by an absolute constant. We also use the ordinary uncentered strong
maximal operator \(M_{\mathrm{str}}\) over axis-parallel rectangles,
which is bounded on \(L^2(\R^2)\).

Let \(Q\subset\mathcal D\) have finite measure, let
\(D\ge D_*\), and put
\[
 a_T(D)=\frac{|Q\cap DR_T|}{|DR_T|}
       =\frac{|Q\cap DR_T|}{D^2|R_T|}.
\]
Here and below dilation is about the box center in its slope
coordinates. We prove, for \(0<\beta\le1\), that
\begin{equation}\label{forest:eq-dilation-tail}
 \sum_{T:a_T(D)>\beta}|R_T|
 \le \min\left(S,\,
       C L^{2C_{\mathrm{as}}}D^2\beta^{-2}|Q|\right).
\end{equation}

Fix a top occurring on the left-hand side. Set
\[
 J_T=I\cap[x_T-Dl_T/2,x_T+Dl_T/2],\qquad a=|J_T|,
 \qquad b=C_{\mathrm{w}}DW,
\]
where \(C_{\mathrm{w}}\ge1\) is a fixed absolute constant.
The witness set is contained in \(DR_T\cap\mathcal D\), so
\(a>0\). The slope parallelogram
\[
 \widetilde R_T=\{(x,y):x\in J_T,\quad
       |(y-y_T)-s_T(x-x_T)|\le b/2\}
\]
contains both \(E_T\) and \(Q\cap DR_T\). It satisfies
\begin{equation}\label{forest:eq-clipped-geometry}
 a\le\min(|I|,Dl_T),\qquad
 ab\le C_{\mathrm{w}}D^2|R_T|,\qquad
 \frac ba\ge C_{\mathrm{w}}\frac W{l_T}
             =C_{\mathrm{w}}m d_T.
\end{equation}
Thus the possibly very long horizontal dilation has been clipped
before any Euclidean enclosing rectangle is constructed.

Suppose first that \(b/a\le\varepsilon_0\), for a fixed small
absolute \(\varepsilon_0\). Projection onto the orthonormal
directions
\[
 e_T=\frac{(1,s_T)}{\sqrt{1+s_T^2}},\qquad
 e_T^\perp=\frac{(-s_T,1)}{\sqrt{1+s_T^2}}
\]
encloses \(\widetilde R_T\) in a Euclidean rectangle \(U_T\)
whose side lengths are
\[
 L_T=\sqrt{1+s_T^2}\,a
           +\frac{|s_T|}{\sqrt{1+s_T^2}}b,
 \qquad B_T=\frac{b}{\sqrt{1+s_T^2}}.
\]
The slopes are absolutely bounded. Consequently \(L_T\asymp a\),
\(B_T\asymp b\), \(|U_T|\le C D^2|R_T|\), and
\[
 L_T\le C|I|,\qquad \frac{B_T}{L_T}\ge c m d_T.
\]
On the witness set, normalization of the direction gives
\(\norm{v_0-e_T}\le C d_T\). Since \(m\to\infty\), this
error is smaller than any required fixed small multiple of
\(B_T/L_T\), for sufficiently large \(L\). The fixed widening
constant \(C_{\mathrm{w}}\) can absorb the absolute aperture
constants as well. Choose the horizontal localization constant so
that \(C|I|\) meets the length restriction in
\eqref{forest:eq-kakeya-input}. The rectangle \(U_T\) is then
admissible with the common popularity parameter
\[
 \delta_D=c L^{-C_{\mathrm{as}}}D^{-2},
\]
where \(c>0\) is sufficiently small and absolute. Indeed,
\(|E_T|\ge L^{-C_{\mathrm{as}}}|R_T|\), and
\(|U_T|\le CD^2|R_T|\). Moreover,
\[
 \frac1{|U_T|}\int_{U_T}\one_Q
 \ge\frac{|Q\cap DR_T|}{CD^2|R_T|}>c\beta.
\]
Every point of \(E_T\) therefore belongs to
\(\{M_{v_0,\delta_D}\one_Q>c\beta\}\).

If \(b/a>\varepsilon_0\), enclose
\(\widetilde R_T\) in the axis-parallel rectangle with horizontal
side \(J_T\) and vertical side of length \(b+|s_T|a\).
Its area is at most \(C ab\le CD^2|R_T|\). It contains
\(E_T\) and \(Q\cap DR_T\), so every point of \(E_T\)
belongs to \(\{M_{\mathrm{str}}\one_Q>c\beta\}\).
This argument also covers boxes whose transverse side exceeds the
horizontal side.

The sets \(E_T\) are pairwise disjoint. Summing their lower measure
bounds and using the two maximal estimates yields
\begin{align*}
 \sum_{T:a_T(D)>\beta}|R_T|
 &\le L^{C_{\mathrm{as}}}
    \left|\bigcup_{T:a_T(D)>\beta}E_T\right|\\
 &\le C L^{C_{\mathrm{as}}}
       (\delta_D^{-1}+1)\beta^{-2}|Q|\\
 &\le C L^{2C_{\mathrm{as}}}D^2\beta^{-2}|Q|.
\end{align*}
The bound by \(S\) is immediate, proving
\eqref{forest:eq-dilation-tail}.

For positive numbers \(S_0,A_0\), direct integration, split at
\(\min(1,\sqrt{A_0/S_0})\), gives
\[
 \int_0^1\min(S_0,A_0\beta^{-2})\dd\beta
        \le2\sqrt{S_0A_0}.
\]
Layer cake in \eqref{forest:eq-dilation-tail} therefore gives
\begin{align}
 \sum_T|R_T|a_T(D)
 &\le C L^{C_{\mathrm{as}}}D\sqrt{S|Q|},
       \label{forest:eq-dilation-layer}\\
 \sum_T|Q\cap DR_T|
 &\le C L^{C_{\mathrm{as}}}D^3\sqrt{S|Q|}.
       \label{forest:eq-dilated-count}
\end{align}
In the second line the additional factor \(D^2\) is the
dilation-area ratio.

Set \(D_k=2^kD_*\). The explicit formula for \(h_T\) implies
the pointwise bound
\[
 h_T(z)\le C_{D_2}\sum_{k\ge0}2^{-kD_2}\one_{D_kR_T}(z).
\]
To verify it outside \(D_*R_T\), choose the least \(k\ge1\)
such that \(z\in D_kR_T\). Then
\(\dist_T(z)>2^{k-2}D_*\), while \(D_*\ge2K_0\), so
\(h_T(z)\le C_{D_2}2^{-kD_2}\). Inside \(D_*R_T\), use
\(h_T\le C_3\). Applying \eqref{forest:eq-dilated-count},
\begin{equation}\label{forest:eq-testing-count}
 \int_Q n
 \le C L^{C_{\mathrm{as}}}D_*^3\sqrt{S|Q|}
       \sum_{k\ge0}2^{-k(D_2-3)}
 \le C L^{C_{\mathrm{as}}}D_*^3\sqrt{S|Q|}.
\end{equation}
The sum converges because \(D_2>4\). Finally take
\(Q=\{z\in\mathcal D:n(z)>\lambda\}\). This set has finite
measure, since the collection is finite and \(D_2>2\). The
inequality \(\lambda|Q|\le\int_Q n\) proves the asserted weak
\(L^2\) bound. All exponents in this argument are fixed: in
particular the dilation cost before the final summation is exactly
\(D^3\), independently of the number of tops and angular scales.
\end{proof}

\subsection{Parameters and masks}

All logarithms below are natural. Set
\begin{equation}\label{forest:eq-parameters}
 \begin{split}
 p&=L^{1/2},\qquad
 G=2^{\,2\left\lceil L^{0.7}/(2\log2)\right\rceil},\\
 K&=\left\lceil(\log L)^2\right\rceil,
 \qquad B=G^K.
 \end{split}
\end{equation}
Thus \(G\) and \(\sqrt G\) are dyadic powers,
\(\log G\asymp L^{0.7}\), and
\(\log B\asymp L^{0.7}(\log L)^2=o(L^{0.9})\).
In particular,
\begin{equation}\label{forest:eq-scale-separation}
 \frac{B^{100}}m\longrightarrow0\qquad(L\longrightarrow\infty).
\end{equation}

The parameters have different tasks. The large bin number \(m\)
permits fine angular phase separation while keeping changes in the
packet envelopes small. The much smaller time cutoff \(\tau\) pays
for the commutation errors. The scale ratio \(G\) will separate
prediction and averaging lengths; grouping \(K\) angular transitions
into a block gives the ratio \(B\). Their strict asymptotic ordering
is what allows these tasks to coexist.

Let \(A>2\) be a fixed exponent to be chosen sufficiently large
after the unmasked top estimates in \eqref{eq:src-11}. Their fixed
exponents and dilation \(K_1=L^{A_1}\) are independent of \(A\).
The comparison of the original and masked sums uses
\Cref{forest:prop-count} at \(K_1\), so its loss
\(c_{\mathrm{cnt}}(A_1)\) is fixed before \(A\) is chosen.
The count at the \(A\)-dependent \(K_0\) below is used only after
that choice. Constants in the mask estimates may depend on the fixed
value of \(A\).
Put \(r_L=2\lceil L\rceil\), and define, in the coordinates
\eqref{forest:eq-coordinates},
\begin{equation}\label{eq:src-8}
 q_T(z)=\left(
       \operatorname{sinc}\frac{X_T(z)}{L^A}\,
       \operatorname{sinc}\frac{Y_T(z)}{L^A}
       \right)^{r_L},\qquad
 \operatorname{sinc}(a)=
 \begin{cases}\sin(a)/a,&a\ne0,\\1,&a=0.\end{cases}
\end{equation}
The exponent is even, so \(0\le q_T\le1\) and the mask is real.
From now on take
\[
 K_0=L^{A+2}
\]
in \eqref{eq:src-5}. A top is called \emph{significant at \(z\)}
if \(\dist_T(z)\le K_0\). Every significant top contributes at
least one to \(n(z)\).

\begin{lemma}\label{forest:lem-mask}
The mask has Fourier support
\begin{equation}\label{forest:eq-mask-spectrum}
 \supp\widehat q_T\subset
 \left\{(\xi,\eta):
       |\xi+s_T\eta|\le\frac{r_L}{L^A l_T},\quad
       |\eta|\le\frac{r_L}{L^A W}\right\}.
\end{equation}
For sufficiently large \(L\), with all other exponents fixed,
\begin{align}
 |q_T(z)|&\le \exp(-L\log L)
       &&\text{if }\dist_T(z)>K_0,
       \label{forest:eq-mask-outside}\\
 |q_T(z)|h_T(z)^{-1}&\le C
       &&\text{for all }z,
       \label{forest:eq-mask-weight}\\
 |1-q_T(z)|&\le C L^{1-A}
       &&\text{if }\dist_T(z)\le L^{A/2}.
       \label{forest:eq-mask-inside}
\end{align}
In particular multiplication by \(q_T\) adds at most
\(O(L^{1-A})\) times the bandwidths of the top box to a Fourier
support.
\end{lemma}

\begin{proof}
For the Fourier convention \(e^{ix\xi}\), the transform of
\(\operatorname{sinc}(x)\) is supported in \([-1,1]\).
Taking its \(r_L\)-th power convolves that transform \(r_L\)
times, so the transform of
\(\operatorname{sinc}(X/L^A)^{r_L}\) is supported in
\([-r_L/L^A,r_L/L^A]\). The product of the two one-dimensional
functions has this support bound in each normalized frequency
coordinate. The affine coordinate change
\eqref{forest:eq-coordinates} sends these coordinates to
\(l_T(\xi+s_T\eta)\) and \(W\eta\), proving
\eqref{forest:eq-mask-spectrum}. The center translation contributes
only a Fourier phase. The calculation can equivalently be performed
first for compactly supported Fourier distributions and then their
convolutions; the resulting mask is integrable for \(r_L>1\).

Write \(d=\dist_T(z)\). Since at least one of
\(|X_T(z)|,|Y_T(z)|\) is at least \(d/2\), the elementary
inequality \(|\operatorname{sinc}(a)|\le\min(1,|a|^{-1})\)
gives, for \(d>0\),
\begin{equation}\label{forest:eq-mask-pointwise}
 |q_T(z)|\le\min\left(1,\left(\frac{2L^A}{d}\right)^{r_L}\right).
\end{equation}
If \(d>K_0=L^{A+2}\), this is at most
\((2L^{-2})^{r_L}\le\exp(-L\log L)\), proving
\eqref{forest:eq-mask-outside}. To include arbitrary distances in
the weighted bound, write \(d=tK_0\), with \(t>1\). Once
\(r_L\ge D_2\),
\[
 |q_T(z)|h_T(z)^{-1}
 \le 2^{-D_2}(1+t)^{D_2}
       (2L^{-2})^{r_L}t^{-r_L}
 \le (2L^{-2})^{r_L}.
\]
For \(d\le K_0\), use \(q_T\le1\) and \(h_T\ge1\).
This proves \eqref{forest:eq-mask-weight}, including the region
arbitrarily far from the top.

Finally, if \(d\le L^{A/2}\), both arguments of the sinc
functions have absolute value at most \(L^{-A/2}\). They lie
in \([0,1]\) after evaluation by sinc, and
\[
 0\le1-
   \operatorname{sinc}(X_T/L^A)
   \operatorname{sinc}(Y_T/L^A)
 \le C\frac{X_T^2+Y_T^2}{L^{2A}}
 \le C L^{-A}.
\]
The inequality \(1-a^{r_L}\le r_L(1-a)\), for
\(0\le a\le1\), proves \eqref{forest:eq-mask-inside}.
Since \(r_L/L^A=O(L^{1-A})\), the bandwidth assertion follows
from \eqref{forest:eq-mask-spectrum} and addition of Fourier
supports under multiplication.
\end{proof}

For completeness, the spectral enlargement also respects the finer
tile coordinates. If \(T=T(s)\), an added frequency
\((\zeta,\nu)\in\supp\widehat q_T\) satisfies
\[
 |\zeta+\theta_s\nu|
 \le\frac{r_L}{L^A}
       \left(\frac1{l_T}+\frac{|\theta_s-s_T|}{W}\right)
 \le\frac{C L^{1-A}}{\ell_s},\qquad
 |\nu|\le\frac{C L^{1-A}}{W}.
\]
Here we used \(\ell_s\le l_T\),
\(|\theta_s-s_T|\le Cd_s\), and
\(d_s/W=1/(m\ell_s)\). Thus all fixed spectral margins needed
later remain available after taking \(L\) sufficiently large.
